\addtocontents{toc}{\protect\setcounter{tocdepth}{-1}}

\begin{essay}{гибридные силовые установки, гибридные трансмиссии, DHT, HEV, MHEV, FEV, PHEV, REEV, Hybrid Synergy Drive, THS, Parallel Inline, Multi-mode Series, топливная экономичность, расход топлива}
Целью данной работы является исследование гибридных силовых установок легковых автомобилей, анализ их архитектурных особенностей и оценка эффективности с точки зрения топливной экономичности и динамических характеристик.

В ходе выполнения работы была изучена современная терминология в области гибридных силовых установок, рассмотрены принципы работы последовательных, параллельных и комбинированных схем гибридизации, а также проведена классификация гибридов по степени интеграции электрического привода. Были проанализированы основные архитектуры гибридных трансмиссий, включая Hybrid Synergy Drive, Parallel Inline, Multi-mode Series и Mild Hybrid, выявлены их конструктивные особенности и области применения. На основе собранных данных по расходу топлива проведён сравнительный анализ эффективности гибридных автомобилей и их аналогов с двигателями внутреннего сгорания в городском, шоссейном и комбинированном режимах движения.

Полученные результаты показали, что гибридные силовые установки обеспечивают значительное снижение расхода топлива, особенно в городском цикле движения. Наибольшую топливную экономичность демонстрируют архитектуры Hybrid Synergy Drive и Multi-mode Series, тогда как Parallel Inline характеризуется стабильной эффективностью на шоссе. Mild Hybrid системы обеспечивают умеренное снижение расхода топлива, но способствуют улучшению динамических характеристик автомобиля. Результаты исследования подтверждают, что выбор архитектуры гибридной трансмиссии должен учитывать условия эксплуатации и требования к балансу между экономичностью и динамикой.

\end{essay}

\addtocontents{toc}{\protect\setcounter{tocdepth}{2}}